\subsection{The Enforcement Agent: Spatial Dynamics and Punitive Actuation}

Operating under the directive of both the local \texttt{BarangayAgent} and the municipal Mayor, the \texttt{EnforcementAgent} serves as the punitive actuator of the simulation. Computationally, this agent is the physical manifestation of the $F$ (Fine) and $P_{\mathrm{Detection}}$ variables found within the Household's utility function.

\paragraph{Part 1: Jurisdictional Typology and Temporal Constraints}

To accurately reflect the decentralized governance of Philippine Solid Waste Management \citep{Nishimura2022}, \texttt{EnforcementAgent} objects are dynamically instantiated under two distinct operational profiles, each with unique computational constraints:

\begin{enumerate}
    \item \textbf{Micro-Enforcement (Barangay Tanods):} Financed by the baseline barangay budget, these persistent agents simulate local village watchmen. Because real-world Tanods are burdened with multiple community duties beyond waste management, their \texttt{step()} function executes a stochastic distraction filter. At the start of each daily tick, a pseudorandom float determines their availability; computationally, there is a 90\% probability (`rand() > 0.10`) that the agent yields its execution thread to other assumed duties and skips garbage monitoring entirely. Their daily apprehension capacity is strictly bottlenecked to a maximum of one non-compliant household.
    
    \item \textbf{Macro-Enforcement (LGU Strike Force):} When the Deep Reinforcement Learning (DRL) Mayor allocates municipal funds to enforcement, it spawns dedicated LGU Inspectors. These agents bypass the stochastic distraction filter, executing targeted patrols every day. They possess a high daily apprehension capacity (up to 25 households) and issue severe municipal citations ($F = \text{\textpeso}1,000$). To prevent memory leaks and simulate the logistical realities of temporary task forces, these objects are bound by a strict 30-day temporal contract. The agent decrements a \texttt{contract\_days} integer daily; upon reaching zero, the object executes a self-destruct sequence (\texttt{model.grid.remove\_agent(self)}), forcing the AI Mayor to continuously justify their budgetary upkeep \citep{TianReview2024}.
\end{enumerate}

\paragraph{Part 2: Spatial Computing and Bounded Rationality}

The physical environment of the Agent-Based Model (ABM) is constructed as a two-dimensional discrete spatial grid. Unlike centralized, equation-based mathematical models that assume uniform mixing or omniscient actors, this spatial formulation enforces strict topographical limitations \citep{Meng2018}. 

Computationally, Enforcement Agents do not possess an $O(1)$ global lookup table of all non-compliant citizens. Instead, they operate under bounded rationality. They must execute localized spatial queries within a defined line-of-sight (e.g., a Moore neighborhood of $Radius = 10$). This architectural design accurately simulates the limited monitoring capacity of real-world enforcers who can only penalize violations physically observed within their immediate vicinity \citep{Dalugdog2021}.

The algorithmic movement of the enforcer relies on a finite state machine bifurcated into two behavioral states:
\begin{itemize}
    \item \textbf{State $S_0$ (Patrol Mode):} The agent conducts a stochastic spatial random-walk, moving to adjacent coordinates. This represents the routine, unguided neighborhood rounds conducted when no immediate violations are visible.
    \item \textbf{State $S_1$ (Chase Mode):} If the $O(r^2)$ spatial query detects a non-compliant \texttt{HouseholdAgent} within the scanning radius, the enforcer dynamically calculates a directed spatial vector toward the closest violator to execute an intervention \citep{Zhao2022}.
\end{itemize}

\paragraph{Part 3: Algorithmic Execution of Punitive Patrols}

During each simulation step, the deployed agent employs this randomized monitoring strategy, selecting a discrete sample of households up to its assigned capacity limit. If a household is flagged as non-compliant, the enforcer invokes the household's \texttt{get\_fined()} method, triggering the psychological and economic mutations detailed in Section 3.3.1. 

Algorithm \ref{alg:enforcement_patrol_english} details the step-by-step heuristic process the enforcement agent undergoes during its daily patrol sequence.

\vspace{0.5cm} 
\begin{algorithm}[H]
\caption{The Enforcer's Daily Computational Patrol Process}
\label{alg:enforcement_patrol_english}
\begin{algorithmic}[1]
\State \textbf{Starting Information Needed:} 
\State \parbox[t]{0.9\linewidth}{$\rightarrow$ The Enforcer's current grid coordinate, a spatial ``Field of Vision'' (radius = 10 blocks), and their daily quota capacity.\strut}
\Statex
\State \textbf{Step 1: Memory Management (Preventing Infinite Loops)}
\State \parbox[t]{0.9\linewidth}{$\rightarrow$ Add the immediately surrounding households to a ``Visited'' memory set. This prevents the computational logic from freezing or getting stuck inspecting the exact same houses on consecutive ticks.\strut}
\Statex
\State \textbf{Step 2: Spatial Query and Movement Logic}
\State \parbox[t]{0.9\linewidth}{$\rightarrow$ Scan the entire 10-block Field of Vision for unvisited \texttt{HouseholdAgents}.\strut}
\State \parbox[t]{0.9\linewidth}{$\rightarrow$ Filter the queried households to isolate those whose current state is ``Non-Compliant.''\strut}
\Statex
\If{there is at least one Non-Compliant household in sight}
    \State \parbox[t]{0.88\linewidth}{$\rightarrow$ \textit{CHASE MODE:} Calculate the shortest path to the violator closest to the Enforcer's current coordinate.\strut}
    \State \parbox[t]{0.88\linewidth}{$\rightarrow$ Move the Enforcer one step along that directed spatial trajectory.\strut}
\Else
    \State \parbox[t]{0.88\linewidth}{$\rightarrow$ \textit{PATROL MODE:} If no active violations are visible, generate a random legal coordinate and take one step in that direction to continue exploring the grid.\strut}
\EndIf
\Statex
\State \textbf{Step 3: Enforcement Execution (The Punishment Protocol)}
\State \parbox[t]{0.9\linewidth}{$\rightarrow$ Establish a tight ``Catch Zone'' around the Enforcer's new location.\strut}
\For{every household physically located inside the Catch Zone}
    \If{the household is actively Non-Compliant}
        \State \parbox[t]{0.85\linewidth}{$\rightarrow$ Trigger the household's penalty function (decreasing their utility, lowering their attitude, and forcing temporary compliance).\strut}
        \State \parbox[t]{0.85\linewidth}{$\rightarrow$ Log the financial penalty into the municipality's overall collected fines ledger.\strut}
    \EndIf
\EndFor
\end{algorithmic}
\end{algorithm}

\paragraph{Part 4: Systemic Feedback and the Probability of Detection}

Crucially, the physical presence of these agents generates a systemic macro-level feedback loop. While the enforcers act strictly locally, their aggregate presence dynamically alters the environment. Following the daily deployment phase, the simulation calculates the barangay's macro-level Perceived Probability of Detection ($P_{\mathrm{Detection}}$) by taking the ratio of the total inspection capacity ($C_{\mathrm{Quarterly}}$) of all currently active enforcers against the total household population:

\begin{equation}
    P_{\mathrm{Detection}} = \min\left(1.0, \frac{\sum C_{\mathrm{Quarterly}}}{N_{\mathrm{Households}}}\right)
    \label{eq:prob_detection}
\end{equation}
\addequation{Macro-level Perceived Probability of Detection}

This updated scalar is broadcasted back to the \texttt{HouseholdAgent} objects, dynamically increasing the anticipated financial risk within their utility calculations. 

This spatial limitation profoundly reflects the logistical realities of municipal governance. Local Government Units (LGUs) rarely have the manpower to achieve $P_{\mathrm{Detection}} = 1.0$ \citep{Abordo2025}. Consequently, the probability of a household being caught is not absolute; it is a highly localized phenomenon dependent on spatial proximity and the enforcer's stochastic patrol route. This mechanism organically generates the dynamic, fluctuating risk profile experienced by citizens, ensuring that the deterrent effect of punitive policies is realistically constrained by the physical capacity of the simulated law enforcement \citep{Jimenez2025}.